\chapter{Single-cell sequencing technologies}
\section{Early single-cell RNA sequencing protocols}
\label{sect:single_cell_sequencing_tech}

% CTCs: Detection, clinical relevance and specific biological properties of disseminating tumour cells

Need for single-cell assays:
\begin{itemize}[noitemsep]
	\item Traditional bulk sequencing techniques mask cellular heterogeneity by averaging across large populations. eg somatic mosaicism in brain neurons,
	\item This is especially important in cancers where rare, therapy-resistant cell types can occur. [Immunology also benefits a ton, TCR sequencing etc. Immune system is very heterogeneous]
	\item Differentiation stages co-existing within a tissue during development
	\item For some samples, most notably early developmental stages, very little material is available.
	\item Capturing expression or accessibility dynamics and noise at the single-cell level
	\item Novel cell type discovery
\end{itemize}


\sjcitep[.]{shapiro2013single}

Curiosity already drove early research to explore gene expression at the single-cell level in the 90s. Eberwine and colleagues injected an oligo-dT T7 primer and T7 reverse transcriptase directly into single rat hippocampal neurons to reverse transcribe their mRNA in situ\sjcitep[.]{eberwine1992analysis} The resulting cDNA was amplified using in vitro transcription, and relative expression levels of 10 known neuronal genes (mostly ion channels) were quantified by Southern blot. Even among \textit{morphologically similar neurons}, significant differences in gene expression were observed, recording for the first time transcriptomic heterogeneity within the same tissue. About a decade later, a double T7-based amplification method was applied to profile some (12) laser microdissected rat hippocampal neurons on cDNA microarrays, revealing two new marker genes for a novel pyramidal neuron subtype\sjcitep[.]{kamme2003singlecell} Soon, PCR-based cDNA amplification strategies\sjcitep[]{brady1993construction,brady2000expression} rather than IVT was applied to single cells from mouse embryonic stem cells and blastomeres\sjcitep[.]{kurimoto2006improved} Mature mRNA was extracted by poly(dT) primers, reverse-transcribed and PCR-amplified, and the resulting single-cell libraries were each hybridised to \textit{individual} RNA microarrays. This revealed that the blastocyst's inner cell mass contains at least two populations: some express genes typical of primitive endoderm (Gata4, Gata6) and others express epiblast markers (Fgf4, Nanog, Klf2).

%Radiolabeled aRNAs were also hybridised to a cDNA library extracted from NG108, a hybrid neuroblastoma-glioma cell line known to express a wide variety of neuron-associated genes. Hybridising clones were isolated and sequenced, leading to the discovery of three novel genes that did not match existing gene databases.

"True" (whole transcriptome) single-cell sequencing was first demonstrated in 2009 by Tang and Surani at the Gurdon Institute\sjcitep[.]{tangMRNASeqWholeTranscriptome2009} By increasing reverse transcription and PCR extension times/cycles, the group was able to amplify enough material for sequencing cDNA libraries from 6 individual mouse blastomeres. This sequencing-based approach provided a number of immediate advantages over microarray analysis:

\begin{enumerate}
	\item Comparison of single-blastomere mRNA-Seq data with pooled (320 cells) Affymetrix microarray data showed that \SI{94.3}{\percent} of microarray-detected genes had at least 5 transcripts in the mRNA-Seq data, but mRNA-Seq detected \num{5270} (up to 75\%) more genes in total.
	\item Sensitivity and dynamic range (\SIrange{0.2}{12000}{RPKM}) were decidedly higher for mRNA-Seq than for analog microarrays.
	\item Since near-whole transcripts were sequenced novel splice junctions could be determined. The group discovered \num{1753} novel isoforms, demonstrated for the first time that different isoforms of the same gene can be expressed in the same single cell.
\end{enumerate}

Tang et al. used the ABI SOLiD sequencer, which had become commercially available only two years prior. Since the sample extraction and amplification protocol itself were only slightly different from existing microarray protocols, it can be posed that development of the Tang method was an effect of sequencing hardware improvements. However, the new sequencing-based approach also came with a new set of challenges. Cost of the novel sequencing technology was very high - each cell in the Tang paper, sequenced at around 24 to 85 million reads, must have cost around \USD{10000} to \USD{50000} in SOLiD sequencing alone \textit{per cell}. At the time, one Affymetric microarray sample cost in the order of \USD{800} - significantly cheaper. Despite longer extension times, the protocol was also unable to capture transcripts longer than \SI{3}{kb}, and was limited to poly(A)\textsuperscript{+} mRNA. These factors contributed to mRNA-Seq's status as a proof-of-concept technique rather than an immediate research tool.

One of the first protocols to be published after the Tang protocol was STRT-seq (single-cell tagged reverse transcription), developed by the Linnarson lab at Karolinska Institute\sjcitep[.]{islamCharacterizationSinglecellTranscriptional2011, islamHighlyMultiplexedStrandSpecific2012} The protocol applied template-switching reverse transcription (use of MMLV reverse transcriptase, which tags 5' transcripts via template switching)\sjcitep[]{zhu2001reverse} to cells picked into individual wells, followed by pooled single-primer PCR for cDNA amplification. A total of 92 mouse cells were profiled: 48 embryonic stem cells and 44 embryonic feeder cells. Libraries were sequenced on the Illumina Genome Analyzer IIx, for a total of 110M raw reads. mRNA capture efficiency was approximately 10\%, as measured by RNA spike-in panels (ERCC). 80\% of filtered reads aligned to exonic and splice-junction features in the mouse genome, which is high even by today's standards. Introduction of degenerate bases (\verb|NNNNNNNNNN|) in the 5' adapter allowed multiple counts of PCR duplicates to be collapsed into transcript counts\sjcitep[.]{kiviojaCountingAbsoluteNumbers2011} Addition of these unique molecular adapters (UMIs) to the STRT-seq protocol, plus parallel automated cell capture on microfluidic arrays (Fluidigm C1) and oligo design optimisation improved the capture efficiency of STRT-seq from \textasciitilde{10}\% to \textasciitilde{50}\% \sjcitep[.]{islamQuantitativeSingleCell2014}

% Later on, STRT-seq was also applied to custom-milled 9600-well nano arrays, recovering between 3000 to 7500 single cells per plate and an average of \num{8706} UMIs and \num{3686} genes\sjcitep[.]{hochgerner2017strt}

% Tumor heterogeneity was already known to underlie treatment resistance\sjcitep[.]{park2010cellular} Bulk DNA sequencing and comparative genomic hybridization assays cannot identify heterogenity in mutation or copy number variation profile. Navin et al. FAC sorted 200 single cells from triple-negative breast cancer tumors, and performed whole genome amplification and sequencing on individual cells\sjcitep[.]{navin2011tumour} Due to low per-cell coverage (\textasciitilde{6\%}), copy number variations was the focus of this study. Lineage analysis identified several "bursts" of CNV shared across main subpopulations, challenging previous ideas of slow, gradual accumulation of mutations underlying cancer. Metastases were strongly enriched for aneuploid populations that represent minorities in the primary tumor, indicating that chromosomal duplication helps establish rapidly expanding metastatic clones. "By isolating, sorting and sequencing the genome of 52 cells from a primary breast tumour and 48 cells from its associated liver metastasis, Navin et al. showed that a single cell from the primary tumor mass seeded the metastatic tumour\sjcitep[.]{navinTumourEvolutionInferred2011}"

In parallel research, also at the Karolinska Institute, Ramsk\"{o}ld and Sandberg developed Smart-seq (switching mechanism at the 5' end of RNA template)\sjcitep[.]{ramskoldFulllengthMRNASeqSinglecell2012} Similarly to STRT-seq, a combination of MMLV reverse-transcriptase followed by cDNA PCR was used. However, fragmentation of these full-length transcripts using Tn5 (or Covaris shearing) allowed for shotgun sequencing of the entire transcript. 42 cells from several human cancer cell lines, circulating tumor cells, and mouse oocytes were profiled. Improved isoform coverage showed that exon 13 of the NEDD4L gene, a prognosis marker, was differentially expressed between LNCaP prostate cancer and T24 bladder cancer cells. A year later, an updated Smart-seq protocol with increased sensitivity was published\sjcitep[.]{picellismartseq2sensitivefull2013} Here, the addition of a locked nucleic acid (LNA) at the TSO's 3' end and addition of betaine led to 2-4 fold cDNA yield increases. The workflow was also simplified and reduced in cost by switching to KAPA HiFi polymerase, which also reduced GC-biases. The use of in-house produced recombinant Tn5\sjcitep[]{picellitn5transposasetagmentation2014} further reduced costs for Smart-seq2\sjcitep[]{picellifulllengthrnaseqsingle2014} - researchers could now bypass Illumina's Nextera Tn5, which to this day is still enormously expensive (\textasciitilde\$\num{3600} per \SI{160}{\micro\liter}). A quantitative comparison of Smart-seq and bulk RNA-seq revealed that aggregate Smart-seq data closely resembles bulk RNA-seq data from 1 million cells, and correlates strongly (Pearson R = 0.9) to ERCC spike-ins\sjcitep[.]{wuquantitativeassessmentsinglecell2014} Moreover, nanoliter volume reactions used in the Fluidigm C1 consistently identified more genes and had more uniform transcript coverage than microliter volume reactions due to higher enzyme, template and substrate concentrations and lower loss in washes. Continued development led to further sensitivity improvements in Smart-seq3\sjcitep[.]{hagemannjensensinglecellrnacounting2020}

CEL-seq, developed in 2012, elected to use linear amplification by in-vitro transcription after reverse transcription rather than PCR-based approaches - effectively translating the Eberwine protocol to sequencing\sjcitep[.]{hashimshonycelseqsinglecellrnaseq2012} The protocol generated \num{280000} mRNA counts from single \textit{C. elegans} blastomeres. An updated protocol later increased low capture efficiency from \textasciitilde{6}\% to \textasciitilde{22}\% and increased throughput to 96-format using the Fluidigm C1\sjcitep[.]{hashimshony2016celseq2}

% This priming site allowed PCR amplification of the resulting cDNA, increasing sensitivity and reducing 3' bias. Second,  Whereas Tang et al. had sequenced just 6 cells, Ramsk\"{o}ld et al. sequenced 42 cells total .  another significant change was that ramsk\"{o}ld et al. had access to significantly cheaper sequencing technology.

% Early Smart-seq studies quickly confirmed significant cellular heterogeneity in tumors. For example, Analysis of \num{430} fresh single glioblastoma cells with Smart-seq showed heterogeneous expression of several receptor tyrosine kinases (EGFR, PDGFRA, FGFR1), potentially leading to redundancy and therapeutic resistance\sjcitep[.]{patelSingleCellRNAseq2014} The existence of CD133+, therapy-resistant "glioblastoma stem-like cells" had been shown before\sjcitep[,]{baogliomastemcells2006}, but single-cell technologies made deeper characterisation possible.

In 2013, Sasagawa et al. argued that Smart-seq, CEL-seq and STRT-seq suffered from amplification biases caused by 3' enrichment, template switching, and pseudo-random Tn5 fragmentation, resulting in poor correlation between technical replicates (Pearson R between 0.7 and 0.85)\sjcitep[.]{sasagawa2013quartz} Their method, Quartz-seq, showed several sensitivity and specificity improvements: by employing Exo I digestion after reverse transcription and cDNA Covaris instead of Tn5 fragmentation, adapter sequences present in the data was reduced from 21\% to 7\%, and technical reproducibility was increased (Pearson R of \textasciitilde{0.93}). Later, Quartz-seq2 would improve throughput by using FAC sorting into 1536-wells, and reduce costs from \$23 to \$0.50 per cell by earlier pooling of the reactions\sjcitep[.]{sasagawa2018quartzseq2} Stringent live-cell selection with FACS, addition of a UMI to the RT primer, optimized poly(dA) tailing of cDNA, and use of Y-shaped Illumina adapters further reduced biases and increased sensitivity, achieving a "UMI conversion rate" of 30-50\%, compared to Quartz-seq1 and Drop-seq's 10\%, and CEL-seq's 22\%\sjcitep[.]{ziegenhainComparativeAnalysis2017, mereu2020benchmarking}

At this point in time, the first plate-based single-cell protocols - particularly Smart-seq - began to spread to other labs, and single-cell sequencing became Nature method of the year in 2013\sjcitep[.]{MethodYear20132014} However, since reagent and labour costs per cell were still high (approximately \$ 30 per cell plus cost of sequencing), almost all studies from this period were limited to few (hundreds) of cells, mainly sourced from few-cell stage embryonic samples\sjcitep[.]{kolodziejczyktechnologybiologysinglecell2015} Profiling samples with inherently large cell counts and diversity, such as immune systems and neuron clusters, would require a different approach.

One strategy to increase cell throughput of pure plate-based methods is to decrease costs by pooling reactions early. In MARS-seq\sjcitep[,]{jaitinMassivelyParallelSingleCell2014} cells are FAC sorted into 384 well-plates, where they undergo cell-barcoded reverse transcription. Reactions are then pooled, and second-strand synthesis and further reactions occur in a shared microtube. cDNA is then amplified by IVT, fragmented, reverse transcribed, amplified via PCR, and sequenced. Applying this to \textasciitilde{\num{5500}} mouse spleen cells led to the discovery of novel dendritic cell subpopulations, exhibiting activation of different Irf4, TNF, and TGF-\textbeta pathways in response to LPS stimulation.

A paradigmatic breakthrough came in 2015, the "annus mirabilis" for droplet-microfluidic single-cell sequencing methods. Researchers at David Weitz' lab had been working on using droplet microfluidics\sjcitep[]{thorsen2001dynamic} to perform antibody screening\sjcitep[]{mazutis2013singlecell} and rare template PCR, referring to a hypothetical system of encapsulating molecules with a nucleic acid barcode that \textit{"indicates which element is encapsulated in the droplet"}\sjcitep[.]{guoDropletMicrofluidicsHighthroughput2012} Such a system was realised for single-cell RNA-seq in 2015 by two separate labs, both at Harvard in collaboration with Weitz lab: Macosko and colleagues (McCarroll lab) developed Drop-seq\sjcitep[,]{macoskohighlyparallelgenomewide2015} and Klein and colleagues (Kirschner lab) developed inDrop\sjcitep[.]{kleinDropletBarcodingSinglecell2015,zilionissinglecellbarcodingsequencing2017}

Both methods use droplet microfluidics to encapsulate single-cells in nanoliter volume droplets together with a bead carrying barcoded poly(dT) primers, but their amplification and library preparation strategies differ. Drop-seq used a Smart-seq-like approach: once encapsulated, cells are lysed, and their polyA\textsuperscript(+) mRNA is captured by the beads, which carry a poly(dT) oligonucleotide primer containing a UMI and a cell barcode. The emulsion is then broken, and beads are pooled and washed in saline sodium citrate to minimize cross-bead hybridizations. The beads, still with mRNA hybridized, are resuspended in a reverse transcriptase mix, and RT proceeds on-bead to produce full-length cDNA. After RT, exonuclease I is applied to degrade any unextended primers, reducing noise from unused barcode sequences. The resulting cDNA, which is covalently linked to the bead, is then PCR-amplified. Fragmentation proceeds using Nextera XT Tn5 to generate a sequencing library. All steps involving enzymes happen in a pooled manner, saving reagent cost. Drop-seq's cell purity (absence of cell doublets) was assessed using a mixture of HEK (human) and 3T3 (mouse) cell lines, and ranged between 90.4\% to 98.8\% based on cell concentration used. The method was demonstrated on \num{44808} cells from the P14 mouse retina generated in 7 runs across 4 days, producing 39 cell clusters. PPP1R17, a new marker for the poorly-characterized nGnG (non-glycinergic/non-GABAergic) amacrine cells, was identified and validated by histochemistry. Transcript capture efficiency was estimated at \textasciitilde10.7-12.8\% using ddPCR validation. Drop-seq marked the first single-cell dataset spanning tens of thousands of cells, and reported a >100-fold reduction in cost and time compared to prior scRNA-seq methods (~6.5 cents per cell, \num{10000} cells in 12 hours). Later, a microfluidics-free version of the protocol, using bead arrays rather than droplets, was developed\sjcitep[.]{gierahnseqwellportablelowcost2017}

inDrop used a different approach to both cell loading and reaction chemistry: cells are encapsulated with soft, deformable hydrogels rather than solid resin beads like Drop-seq, which can be stacked in the microfluidic channel. This ordered stack of beads allows almost every droplet to receive a bead, rather than  >\textasciitilde{2}\% for the poissonian loading used in Drop-seq\sjcitep[.]{abate2009beating} Reaction chemistry was also different. After encapsulation, UV light releases the barcoded poly(dT) primers from the bead into the droplet, and reverse transcription happens within the droplet. Next, the CEL-seq IVT strategy is used to amplify the barcoded, single-cell transcriptomes in a pooled manner. The protocol was demonstrated on a total of \num{11209} mouse embryonic stem cells, including a time series of leukemia inhibitory factor (LIF) withdrawal-induced differentiation over 7 days. 5 distinct ES sub-populations were identified, including rare lineages marked by epiblast (Krt8/18), primitive endoderm (Col4a1/2, Lama1), and heat shock markers (Hsp90, Hspa5). Counts agreed strongly with smFISH data (R = 0.92 for mean, 0.90 for CV).

Both Drop-seq and inDrop boasted a significant jump in throughput over pure plate-based sorting methods, from hundreds to several thousands, potentially tens of thousands, of cells per day. Higher cell counts, traded for per-cell sensitivity, now allowed for the profiling of highly varied samples such as immune cells. However, both methods also required fabrication of complex reagents (notably, barcoded beads and microfluidic chips) and complex microfluidic operation. In 2016, the first commercial droplet-microfluidic single-cell sequencing solution was released: 10x Genomics' 3' Single Cell Gene Expession kit and accompanying Chromium microfluidic controller combined barcoded hydrogel beads, enzymes, primers and microfluidics in a relatively simple package. The method drew on concepts from both inDrop and Drop-seq, opting for deformable hydrogel beads for higher cell capture, and a SMART-seq-like, PCR + Tn5 amplification and library preparation strategy. Importantly, the Chromium could load up to 8 samples in parallel, simplifying arrayed experiments. The method was first demonstrated on \num{68000} PBMCs sampled from 8 donors, identifying all major PBMC subtypes (T, NK, B, myeloid) and minor populations (e.g. megakaryocytes)\sjcitep[.]{zhengmassivelyparalleldigital2017} Approximately \num{4500} genes and \num{27000} UMIs were detected on average per cell at a depth of 100k reads/cell, an order of magnitude higher than inDrop and Drop-seq. Mixed-species loading revealied a doublet rate of approximately 3\%.

The release of the 10x Chromium scRNA-seq kit marked a milestone in single-cell sequencing research. It was user friendly enough, and struck a good balance between cell throughput and single-cell quality. Several atlasing projects were completed using the technology, including human\sjcitep[,]{regev2017human} mouse\sjcitep[,]{tabulamuris2018} and fly\sjcitep[.]{davie2018single} In the latter, \num{157000} high-quality single-cell \textit{D. melanogaster} transcriptomes were generated in 26 runs sampled across 8 timepoints between 0 and 50 days old. CEL-seq2 and SMART-seq2 supplemented the medium-depth 10x Chromium cells with high-quality "anchors" of higher sequencing depth. New, cell-type-specific gene regulatory networks were inferred from cell-type-specific (olfactory projection neurons, glia) gene sets using SCENIC. Participant TFs were validated for motif enrichment in these cell types using FACS-assisted ATAC-seq.

% 10x's Chromium drew on some technical concepts borrowed from droplet digital PCR (ddPCR), such as emulsion chemistry and microfluidic geometry, which had previously been developed by the company's CSO at Lawrence Livermore\sjcitep[,]{hindsonHighThroughputDropletDigital2011} and commercialized at QuantaLife\sjcitep[.]{brittain2021biorad}

% 2016: SCRB-seq\sjcitep[]{soumillonCharacterizationDirectedDifferentiation2014}
% 2018: mcsCRB-seq\sjcitep[]{bagnoli2018sensitive}

Another independent, mechanistically different, leap in cell-throughput was made in 2015 when researchers in Shendure lab applied combinatorial indexing to single cells in-situ\sjcitep[.]{cusanovich2015multiplex} In sci-ATAC (single-cell combinatorial indexing ATAC), nuclei are tagmented using Tn5 loaded with an array of barcoded adapters in 96-well plates and re-distributed into a second 96-well plate containing PCR mixes with 96 different primer combinations. Each nucleus is thus indexed by a combination of its Tn5 and PCR barcode, and individual cells are never physically separated. The method was applied to \num{14533} cells from 4 different cell lines. Between 50 and 60\% of reads were located in ENCODE-defined DNaseI hypersensitivity sites, confirming assay fidelity. Cross-cell collision rate was estimated at 11\% using species identities in the cell lines, and individual cell barcodes contained \textasciitilde90\% species-pure accessible fragments. Interestingly, 4 different states were observed in GM12878 cells based on accessible enhancers. The same combinatorial indexing strategy was applied to scRNA-seq in 2017, this time profiling \num{49360} \textit{C. elegans} L2 larvae, approximately 50-fold coverage over its 762 somatic cells at this stage\sjcitep[.]{caoComprehensiveSinglecellTranscriptional2017} Instead of $96 \times 96$, split-pooling was performed over $576 \times 960$ possible wells (first during in-situ reverse transcription, then Tn5 fragmentation and PCR). Median UMIs and genes detected in \textit{C. elegans} was 575 UMIs and 431 respectively, identifying 27 major cell types, including rare neurons with only 1 or 2 cells per larva. Reagent costs were estimated at \USD0.03-\USD0.20 per cell. Updated protocols were published in 2022, offering higher sensitivity and lower costs (\USD0.01 per cell)\sjcitep[.]{martinOptimizedSingleNucleus2023} SPLiT-seq (split-pool ligation-based transcriptome sequencing), an independent combinatorial indexing protocol, was used to generate \num{156049} single-cell transcriptomes from P2 and P11 mouse brains and spinal cords. This technology introduced additional barcode ligation steps to increase barcode complexity to 96\textsuperscript{4}. The protocol was recently commercialized into kits by Parse Biosciences\sjcitep[.]{kuijpers2024split} Plate-based combinatorial indexing strategies have in general skirted on the extreme ends of the sensitivity/cell count trade-off, and have been used in million-cell scale atlasing projects over the years (\autoref{table:millions}).

\begin{table}[h]
	\centering
	\caption{Large-scale single-cell atlases using sci-RNA-seq3}
	\begin{tabular}{
		>{\centering\arraybackslash}m{1.5cm}
		>{\centering\arraybackslash}m{2.2cm}
		>{\centering\arraybackslash}m{3.5cm}
		>{\centering\arraybackslash}m{3.5cm}
		}
		\toprule
		Year                                                                         & Cells                   & Sample                  & Technique    \\
		\midrule
		2019\sjcitep[]{cao2019singlecell}                                            & \num{2072011}           & \makecell{E9.5 to E13.5                \\ mouse embryos} & sci-RNA-seq3 \\
		2020\sjcitep[]{caoHumanCellAtlas2020}                                        & \num{4062980}           & Human embryos           & sci-RNA-seq3 \\
		2023\sjcitep[]{chiousinglecellmultiomic2023}                                 & \makecell{\num{2583967}                                          \\ + \num{1587880}} & Macaque brain & \makecell{sci-RNA-seq3 \\ and snATAC-seq3} \\
		2024\sjcitep[]{martinOptimizedSingleNucleus2023, qiuSingleCellTimeLapse2024} & $12.4 \times 10^6$      & \makecell{E8 to P0                     \\ mouse embryos} & sci-RNA-seq3 \\
		\bottomrule
	\end{tabular}
	\label{table:millions}
\end{table}

ATAC-seq:
chen protocol\sjcitep[]{chen2018rapid}
10x\sjcitep[.]{satpathyMassivelyParallelSinglecell2019}
ddSEQ\sjcitep[.]{lareaudropletbasedcombinatorialindexing2019}
s3atac\sjcitep[.]{mulqueenhighcontentsinglecellcombinatorial2021}

\section{Multimodal technologies}
sci-car\sjcitep[.]{cao2018jointprofiling}

tackling multimodality: not a great strength of droplet microfluidics
historically, most multiomics methods have relied on plates. this makes sense, because you have to do a number of manipulations/washes on the cells in order to perform all the different omics steps. sci-car. nome-seq, voet etc.

droplet methods:
\begin{itemize}[noitemsep]
	\item reap-seq, cite-seq, eccite-seq
	\item 10x multiome\sjcitep[]{schwendinger2020chromium}
	\item TEA-seq\sjcitep[]{swanson2021simultaneous}
	\item SHARE-seq\sjcitep[]{ma2020chromatin}
	\item ASAP-seq\sjcitep[]{mimitou2021scalable}
\end{itemize}

% more complex assays coming up\sjcitep[.]{deJongheSpinDrop2023}
% digital microfluidics\sjcitep[.]{lamannaDigitalMicrofluidicIsolation2020}


\section{Perturbations}
Large screens that investigate the effect of silencing specific genes using RNA interference emerged in the early 2000s\sjcitep[.]{berns2004large,boutros2006analysis} Such perturbation screens follow either an arrayed design, where perturbations are individually plated and tracked in physical space, or a pooled design, where cell lines are transduced with pooled shRNA vector libraries at a given MOI, and survivors are sequenced.

% TODO: check this last line, how were these pooled screens done?

When CRISPR was developed as a molecular tool for rapid and accurate gene editing in 2013\sjcitep[,]{ran2013genome} full genetic knockouts became widely available, producing a pooled genome-wide screen involving 64000 sgRNAs for 18000 gene knockouts just a year later\sjcitep[.]{shalem2014genome, behan2019prioritization, pacini2021integrated} In addition to genetic knockouts, specific genes can also be up- or downregulated using catalytically dead Cas9 fused to transcription activators such as VP64 or p65, or repressors such as KRAB. Both arrayed and pooled CRISPR/RNAi screens have drawbacks: arrayed screens require vast amounts of labour and reagents, and pooled CRISPR-screens are limited to very simple readouts such as cell survival or marker expression.

Around 2015, researchers at the Broad institute devised that perturbations could be pooled and read out at the single-cell level, coupling pooled screens with scRNA-seq readings\sjcitep[.]{dixit2016perturbseq} In their method, named Perturb-seq, each vector is designed to express two transcripts after integrating in the host genome: a first transcript containing the sgRNA, driven by U6, and a second transcript containing a puromycin selection cassette, BFP and a 20 nucleotide long guide barcode (GBC), driven by EF1\textalpha. The GBC transcript is polyadenylated, allowing it to be captured and counted through standard 3' scRNA-seq protocols (10x GEX in this specific case). The \{sgRNA:GBC\} dictionary can then be used to identify the genetic edit present in each captured cell \sjcitep[.]{dixit2016perturbseq} This pooled approach allows rapid parallelisation of perturbation screens with limited labour. Using this system, researchers observed that even within the same perturbation, individual cells can exhibit different responses. In particular, knockdown of the key ER chaperone BiP (HSPA5) caused two possible responses: some cells preferentially activated PERK, while others activated IRE1\textalpha.

\begin{figure}[h]
	\centering
	\includegraphics[width=0.75\textwidth]{perturbation_vectors_cropped.png}
	\captionsetup{width=0.85\linewidth}
	\caption[Common single-cell perturbation vector designs]{Common single-cell perturbation vector designs. Adapted from \notefullcite{xie2018frequent}.}
	\label{fig:perturbation_vectors}
\end{figure}

Instead of guide barcodes expressed on the selection cassette, a poly-adenylated gRNA which is directly captured by 3' scRNA-seq can also be used\sjcitep[.]{datlinger2017pooled}. In this approach's debut, Drop-seq was also used instead of 10x, showing that the open-source method has sufficient sensitivity to detect the gRNA transcripts. Perturb-seq's costs can be further scaled down by multiplexing perturbations or cells: if gene regulatory networks are assumed sparse in number of participants, then the differential expression of a cell responding to multiple perturbations (high MOI), or a droplet containing multiple single-perturbation cells, can be considered a linear combination of the effect of its components. Pooling cells or perturbations in this manner and demultiplexing the individual terms using PCA and LASSO [double check] can reduce cost ten-fold and quantify interaction terms between different perturbations\sjcitep[.]{yao2024scalable}

Perturbation screens have also been coupled with scATAC-seq readouts. In the earliest protocol, Perturb-ATAC\sjcitep[,]{rubin2019coupled} cells were transduced in pools with Perturb-seq-like vectors. Single cells were then captured and tagmented on the Fluidigm C1, limiting the size of the study to \~4300 cells total. Here, the GBCs are reverse-transcribed, allowing them to be amplified and captured through the Fluidigm scATAC-seq protocol, generating perturbation-coupled scATAC-seq profiles. Dual knockdown of EBF1 and EZH2 (in B lympholbasts), and KLF4 and ZNF750 (in keratinocytes), showed stronger repression than individual perturbations, and this relationship is cell-type specific.

Spear-ATAC\sjcitep[]{pierce2021high} is able to overcome Perturb-ATAC's throughput limitations by using droplet encapsulation (10x Chromium). After pooled lentiviral transduction, cells undergo the regular 10x scATAC-seq workflow, but with an additional primer added to the droplet encapsulation step. This primer anneals to the integrated lentiviral cassette, amplifying its barcode. By using droplet-based methods, 414 sgRNA knockdowns could be profiled across 105k cells (sampled between 3 and 21 days post transduction). Analyzing the resulting time series revealed time-dependent effects of the knockdowns. For example, GATA1 knockdown was found to first upregulate STAT5 motif accessibility, but later RUNX motifs. As expected, these GRNs were also observed to be cell-type specific between K562, MCF-7 and GM12878 cells.

One disadvantage of the screening methods above is that most of the vectors described do not contain Cas9, limiting the studies to transgenic and immortal cell lines. Most perturbation studies so far also deal with CRISPRa/i, CRISPR knockouts, but not yet with synthetic gene expression knock-ins. Such studies are considerably more difficult to implement in a pooled manner, since transcript length has an effect on transfection, packaging and transduction efficiency, which bias MOI. These effects call for the titration of viral concentrations and arrayed screening, which is expensive and laborious.

\begin{itemize}[noitemsep]
	\item crop-seq
	\item TAP-seq
	\item MIX-seq
	\item DRUG-seq -> also interplay with automation
\end{itemize}

\section{Multiplexing}
Sample multiplexing for droplet microfluidics has also been achieved through other means. The main strategy is to pre-index or "hash" multiple samples with oligonucleotide barcodes, co-process them through droplet generation together, and use the sample hash to computationally separate cells into sample bins post hoc. A number of different technical approaches to do this exist. Cells or nuclei can be tagged by attaching sample-barcoded oligonucleotides to their membranes, either by antibody-antigen binding\sjcitep[]{stoeckiusCellHashing2018} or lipid/cholesterol embedding\sjcitep[.]{mcginnisMultiSeqSample2019} Within each droplet, the membrane-bound oligonucleotides co-captured together with the associated cell's mRNA or tagmented genome. Antibodies have demonstrated slightly better accuracy, while lipid hashing can be used when high-strength antibodies are not available\sjcitep[.]{mylkacomparativeanalysisantibody2022} The increased precision of antibody-based methods could be attributed to the many successive washes involved, but which may also reduce cell viability and integrity.

Another sample multiplexing strategy is scifi-RNA-seq, where cells or nuclei are distributed over 96-well plates and pre-indexed by reverse transcription with a barcoded primer akin to plate-based combinatorial indexing protocols. Cells can then be pooled and run through droplet generation together\sjcitep[.]{datlingerultrahighthroughput2021} Beyond just sample multiplexing, scifi-RNA-seq offers a major advantage over lipid and antibody hashes: since each cell's nucleic acid species are pre-indexed, multiple cells can be loaded into the same droplet, which can be used to multiply the total number of cells profiled per droplet run. Whereas the 10x Chromium was designed for \textasciitilde\num{10000} cells per channel, scifi-RNA-seq increases this to up to 1.53 million cells per channel without apparent reduction in per-cell data quality. The same principle can be applied to ATAC-seq: scifi-ATAC-seq and txci-ATAC-seq use pre-indexing in 96 wells with barcoded adapter Tn5, enabling the profiling of up to \num{200000} single nuclei per 10x scATAC-seq lane\sjcitep[.]{zhang2024scifi, zhang2024txci} Sensitivity remained high, sampling between \num{21362} and \num{45622} unique fragments per cell on average across tissues assessed.

SNuBar (single-nucleotide barcoding) introduces a single oligonucleotide barcode (SBO) during tagmentation\sjcitep[.]{wang2021simple} This SBO contains an 18 bp sample barcode and hybridises to Tn5 transposomes through a universal adapter sequence. During tagmentation these Tn5-SBO complexes remain stable as Tn5 integrates itself into accessible regions. Nuclei are then washed, removing unbound Tn5 and excess SBOs, and samples are pooled. After droplet encapsulation, the SBO is released from Tn5 and captured by the cell barcode during PCR amplification, linking cell barcodes with sample barcodes.

Overloading and unpacking (OAK)\sjcitep[]{wuOverloadingAndUnpacking2024} applies a different solution to achieve the same result: cells are reverse-transcribed inside droplets on t

\begin{itemize}[noitemsep]
	\item genetic demultiplexing
	\item vector based demultiplexing https://pmc.ncbi.nlm.nih.gov/articles/PMC11073057/
\end{itemize}

% [draft idea: layered approach (like in davie2018): sci, plate, droplet has different strengths, use their strengths and supplement each other? potentially in a decentralized academic context with many participants, in a centralized effort it's probably better to hard focus on one technique and scale it out]